\subsection*{IV 20.8: Dimension de l'ensemble des hyperplans exceptionnels (cote 25, sites 31--54)}
This long, repeatedly revised section studies the codimension of the exceptional hyperplanes for the properties developed in sections 2--7, with linear pencils as its stated motivation. The opening distinguishes the constructible exceptional set from its closure and asks whether these refinements belong here or beside the earlier results; the first version of 8.2 is canceled on site 31. Its replacement on 32 treats excess dimension and then the regularity of the section $\phi_\xi$ on a coherent Module $F$, using associated prime cycles and the incidence family. Here the first assertion literally says codimension2, while the later Module formulation has an explicit comparison sign before2; the first statement is not supplied with one. The same discussion changes among $Z$, $T$, and their irreducible components. Site 33 passes through coprofondeur, the Cohen--Macaulay property, and $(S_k)$: the added sentence near the top has $(S_{k-1})$, whereas the final handwritten continuation has $(S_k)$. The canceled intermediate argument and its replacement must be read with that distinction of written indices in view. Proposition 8.4 and its dimensional lemma, introduced as a remords to paragraph 5, prepare a more precise description of the exceptional hyperplanes and then the criterion for $(S_k)$ in8.5--8.6. The author cites the dimension theory of $0_{\mathrm{IV}}$ and published EGA IV 5.7.4. Several literal forms require attention: site 34 says $T_i$ is a component ``de X'', gives the lemma the inserted hypothesis $\operatorname{dodim}(Y,X)\geq k$, and cancels $\dim f(X_j)$ while leaving the following equals-zero expression; its last line has $f(T)$ without an index. The corollary on 35 is replaced rather than merely augmented: the new introductory sentence cites8.4,5.7.4, and the end of 8.2, and the partly clipped handwritten hypothesis remains uncertain. Site 36 retains competing readings in $n+k+1$, $f(T_j)$, and the $F/G$ expression; its necessity argument literally calls $\phi_\xi$ $X_{k(\xi)}$-regular amid formulas about $F$ and $G_\xi$. The next part, beginning with the revised number 8.6 on 37, gives the differential description behind the exceptional locus. It uses the exact sequences on projective space, the conormal Module, and the singular and supersingular zero constructions of paragraph 16. The projective conormal description and the diagrams on 39--42 connect ordinary singular points, ramification, tangent hyperplanes, and the stated orthogonality. The author's canceled preliminary argument on 39 gives $M$ rank$n$, writes $\underline{O}_X^{\,m}$, and uses $n-m$ for the fiber dimension; those three expressions are retained together. Other written notation is likewise left visible: the $r-1$ and $d-1$ occurrences on 38, the change from $\psi$ on 39 to $\mu$ on 40, the term $\Omega_{X/Y}$ in40's diagram, and the $X/Y$ subscript and exponent$-2$ in42's lower-left object. The hypotheses and numbering also evolve:8.7 changes the point label from$x$ to$y$,8.8 receives the immersion hypothesis and the reminder that$X$ remains smooth, and the margin of 43 restricts the role of immersion to the equivalence with clause(i ter). The dimension sentence on 43 has the handwritten alternative $\underline{Y}^{\mathrm{sup\ sing}}\neq\emptyset$ followed by the revised equality$r-2$; the later inserted alternative on 45 has an equals-empty-set sign and belongs immediately before the dimension statement there. The edition preserves both signs. The proof on 44--45 works with the tangent map and its image and with a determinant section; the first handwritten morphism label on 44 remains uncertain, and 45 retains both the earlier$r$ and the later$r-1$ dimension expressions. Sites46--49 return to regularity, geometric normality, irreducibility, and the extent of the required nonramification hypotheses. The term ``réglé'' and successive versions of the exceptional-set criterion are central here. Long canceled formulations on 47--48 remain in the diplomatic text and apparatus; the repeated call marks on 47 join the top addition to the sentence about what$T$ contains, while those on 48 place the addition beginning ``grâce au fait'' after the codimension bound. That latter addition visibly uses $\underline{X}'^{\mathrm{sing}}$, and the text retains it. The correction note on 49 keeps its two primed$X'_i$ expressions and a doubtful$X/Y$ in the final image formula. The remords on 50 weakens the proposed nonramification assumption;8.16 and 8.17 then combine geometric normality, depth, connectedness, and irreducibility. The codimension character in the remords remains an unresolved overprint. The quadric example contains the handwritten ``en effet tous'', giving the literal wording about tangent hyperplanes being reducible; on 51 the author immediately records his alarm that the quadric does not qualify under his own use of ``réglé'' and asks for another word. This reservation is part of the argument's unfinished state. Proposition 8.18 on 51 has an inserted dimension$d\geq2$, still says dimension1 in its main assertion, and changes the dimensions in clauses(a) and(b) to$d-1$; clause(a) retains an unchecked$P(k)$. The top handwritten addition, placed before ``Nous allons'' by matching call marks, uses an underlined$Y^{\mathrm{sing}}$ without a$\xi$ index. Its final short mark remains doubtful. The canceled proof below it gives way to 52's proof through the normal case, the open subset$U$, and normalization; the handwritten note compares8.18 with8.16 and explains the interest in allowing$d\geq2$. The closing discussion on 53--54 studies hyperplanes containing a prescribed linear subspace through the projection $P-L^o\to Q$ and invokes the earlier Bertini--Zariski argument of 3.1. Site 53 changes from$Y_\xi$ to$X'_\xi$ and$X_\xi$, then literally ends the equivalence with unindexed$X$; its proposed converse also retains the semicolon and doubtful numeral after$\dim f(X)$. The generic-fiber condition is added below that page and developed in the transcendence-degree discussion on 54. The final dimension estimates, the condition at the finitely many points of$T$, and the suggestion to move these ``Bertiniesques'' developments to the following section close the passage. No index, inequality, hypothesis, or diagram has been repaired from this commentary. Sources: \href{https://grothendieck.umontpellier.fr/web/viewer.html?file=../25.pdf\#page=32}{sites 31--36}, \href{https://grothendieck.umontpellier.fr/web/viewer.html?file=../25.pdf\#page=38}{sites 37--42}, \href{https://grothendieck.umontpellier.fr/web/viewer.html?file=../25.pdf\#page=44}{sites 43--49}, \href{https://grothendieck.umontpellier.fr/web/viewer.html?file=../25.pdf\#page=51}{sites 50--52}, and \href{https://grothendieck.umontpellier.fr/web/viewer.html?file=../25.pdf\#page=54}{sites 53--54}.
